% Generated draft: requires mathematical and source review.
\documentclass[11pt]{article}
\usepackage[margin=25mm]{geometry}
\usepackage{fontspec}
\setmainfont{Latin Modern Roman}
\usepackage{amsmath,amssymb,mathtools,unicode-math}
\setmathfont{Latin Modern Math}
\usepackage{enumitem,needspace,xurl,hyperref}
\setlength{\parindent}{0pt}
\setlength{\parskip}{5pt}
\setlength{\emergencystretch}{3em}
\newcommand{\N}{\mathbb{N}}
\newcommand{\Z}{\mathbb{Z}}
\newcommand{\Q}{\mathbb{Q}}
\newcommand{\R}{\mathbb{R}}
\newcommand{\C}{\mathbb{C}}
\newcommand{\F}{\mathbb{F}}
\newcommand{\A}{\mathbb{A}}
\newcommand{\PP}{\mathbb{P}}
\newcommand{\E}{\mathbb{E}}
\newcommand{\eps}{\varepsilon}
\newcommand{\dd}{\,\mathrm{d}}
\begin{document}
\section*{UnsolvedMath ID: 11000170}
\textbf{Source problem code:} catalogue entry
\subsection{Definitions and mathematical statement}
For a closed oriented surface $S_g$ and homeomorphism $\varphi:S_g\to S_g$, its mapping torus is
\[
 M_\varphi=(S_g\times[0,1])/((x,1)\sim(\varphi(x),0)).
\]
It is a surface bundle over $S^1$; for trivial monodromy it is $S_g\times S^1$. This is the conventional setting suggested by the excerpt, whose exact surface and genus hypotheses are not supplied.

Problem 2.13 reports that [43] proves uniqueness up to equivalence of the bundle-related splitting for $S_g\times S^1$, and asks for other cases in which it is unique. The preceding definition of ``bundle-related splitting'' and its equivalence relation are missing. If these are a specified class $\mathcal B(M_\varphi)$ of Heegaard splittings, the question asks for further bundle cases with
\[
 \#\bigl(\mathcal B(M_\varphi)/\text{the source equivalence}\bigr)=1.
\]
A Heegaard splitting conventionally decomposes a closed $3$-manifold into two handlebodies along their common boundary, but the excerpt does not identify its permitted bundle-related subclass or genus. Uniqueness among that subclass must not be replaced by uniqueness among all splittings, all genera or stabilizations. Nor can equivalence be silently taken as isotopy, homeomorphism, or bundle-preserving equivalence.
\subsection{Short English statement}
Beyond the product surface bundle, when is the source's bundle-related splitting unique up to its specified equivalence? The class of splittings and the equivalence notion are missing.
\subsection{Sources}
\begin{itemize}
\item Local registry: \url{data/unsolvedmath/problems.json}, permanent UnsolvedMath ID $11000170$, exact source code \texttt{catalogue entry}.
\item Farb, \emph{Problems on Mapping Class Groups and Related Topics}, Problem 2.13, PDF page 159: \url{https://www.math.uchicago.edu/~farb/papers/mcgbook.pdf}. Supplied access date: 2026-07-29; extraction: pdf-font-structure; repeated local numbering disambiguated by page among 329 headings.
\item Source reference [43] is cited for the product-bundle uniqueness result, but its bibliography entry and precise theorem are absent. The registry literature summary is truncated.
\end{itemize}
\end{document}
